% Partie IS (source : ../is/note-is.tex, corrigée selon ../revue/verif_is_pfu.md ; chiffres : ../is/sorties/simulation_is.txt, euros de 2027)
L'impôt sur les sociétés (IS) n'appelle pas de grande réforme. Ce que lui reprochent les comparaisons
internationales, c'est son taux affiché, le plus élevé de l'OCDE pour les plus grands groupes en 2025 et 2026
avec la contribution exceptionnelle, sa complexité et son instabilité, pas le traitement de
l'investissement, qui est dans la moyenne. La référence garde donc le taux normal à 25~\% ; la question est
de savoir s'il faut déduire l'investissement plus vite, et comment le payer.

\subsection{Ce que les comparaisons reprochent à l'IS français}

La Tax Foundation classe la France dernière des 38 pays de l'OCDE en 2025 et en 2026, au total comme pour
l'IS, à cause du taux de 36,1~\% des plus grands groupes (24,2~\% en moyenne), de quatre taux différents
(2,3 en moyenne), d'une surtaxe, d'une boîte à brevets et de la troisième subvention implicite à la
recherche la plus généreuse. Cet indice pénalise par construction toute incitation ciblée. Il ne pénalise
pas la déduction de l'investissement : la France est 15\up{e} sur 38 en 2025 et 16\up{e} en 2026 pour la
récupération des coûts, et les déductions d'une machine y valent 88~\% de son prix en valeur actuelle, contre
86~\% en moyenne. Les taux effectifs de l'OCDE vont dans le même sens. Le taux moyen effectif, qui compte pour
choisir où implanter un projet rentable, est de 23,6~\% contre 22,2~\% (13\up{e} sur 38) ; le taux marginal
effectif, qui compte pour décider d'un investissement de plus, est de 15,0~\% contre 16,5~\%, et de 9,8~\%
contre 10,4~\% pour les équipements (chiffres de 2025, qui semblent ignorer la contribution exceptionnelle).
Le Conseil des prélèvements obligatoires (2023, 2025) et le Sénat (2026) font de l'instabilité le premier
grief : la contribution exceptionnelle, présentée comme temporaire, a été prorogée une fois, et son assiette
est la moyenne de l'IS de deux exercices, dont l'un était clos quand elle a été votée.

\subsection{Déduire l'investissement plutôt que baisser le taux}

Une entreprise achète une machine de 100~€ qui dure sept ans. Elle en déduit aujourd'hui le coût par
fractions ; actualisées à 5~\%, ces déductions valent 85 à 90~€ selon que l'amortissement est linéaire ou
dégressif. Si elle peut déduire les 100~€ l'année de l'achat, le total ne change pas, seulement sa date ; mais
l'État rend alors 25~€ à l'achat et prélève ensuite 25~\% de ce que la machine rapporte. Il devient un
associé, à condition que les pertes soient elles aussi déduites sans limite, ce que le plafonnement de leur
imputation empêche en partie. Un investissement qui rapporte tout juste le rendement normal du capital ne
paie plus d'impôt ; restent imposés les profits qui dépassent ce rendement et ceux du capital déjà installé.
Une baisse du taux allège au contraire tous les profits, y compris ceux des machines achetées il y a dix ans et
ceux des rentes, qui ne dépendent pas de la décision d'investir. Sur la réforme américaine de 2017, un point
de capital supplémentaire au bout de dix ans coûte environ 10~\% de la recette de l'IS par an avec une baisse
du taux et 2,5~\% avec la déduction immédiate (Chodorow-Reich, Smith, Zidar et Zwick, 2026). L'effet sur
l'investissement est établi : +10 et +17~\% d'investissement éligible pour les déductions partielles
américaines de 2001-2004 et 2008-2010, deux fois plus dans les petites entreprises et seulement chez celles
qui payaient l'impôt (Zwick et Mahon, 2017) ; +11~\% pour les entreprises moyennes britanniques (Maffini, Xing
et Devereux, 2019). Ces mesures étaient temporaires et ont en partie avancé des achats ; d'où l'importance
d'une mesure permanente. Il n'existe aucune évaluation causale du suramortissement français de 2015.

Le coût d'usage du capital, ce que doit rapporter une machine pour couvrir son financement, son usure et
l'impôt, s'écrit $c = (r+\delta)(1-\tau z)/(1-\tau)$, où $z$ est la valeur actuelle des déductions pour 1~€.
Avec $z$ entre 0,85 et 0,90, l'impôt le majore de 3,5 à 5~\% ; avec la déduction immédiate, plus du tout. Si
le stock d'équipements des sociétés imposables (environ les trois quarts des 541~Md€ des sociétés non
financières) réagit à peu près proportionnellement au coût d'usage (élasticité de 0,5 à 1,6), le PIB gagnerait
0,04 à 0,2~\% à long terme, 0,1~\% au centre. Ce gain de production n'est pas un gain de bien-être du même
montant : il rémunère d'abord le capital supplémentaire qui le produit. La perte sèche évitée, le triangle
sur ce capital, n'est que de 0,05 à 0,16~Md€ par an, parce que le coin supprimé est petit ; l'essentiel de
l'intérêt de la mesure tient à ce qu'elle coûte peu.

\subsection{Combien cela coûte, et comment le payer}

Les sociétés non financières ont investi 84~Md€ en 2025 en machines, équipements et matériels informatiques ;
on suppose que les trois quarts sont faits par des sociétés imposables, hypothèse non mesurée à laquelle le
coût est proportionnel. La première année (2027), une déduction immédiate totale coûterait 14 à 15,5~Md€ ;
le coût baisse ensuite, parce que les machines achetées après la réforme n'ont plus d'amortissement à
déduire. À partir de la huitième année ne reste qu'un coût permanent de 1,4 à 2,0~Md€ par an, qui vient
seulement de ce que l'investissement croît. En valeur actuelle, une année d'investissement coûte 1,8 à
2,6~Md€ au taux d'actualisation des entreprises et 1,2 à 1,7~Md€ à celui de l'État, soit 0,5 à 0,8 point de
taux normal (un point vaut 3,2~Md€ en 2027 une fois le taux réduit supprimé).

\figbloc{is_01_cout_deduction}{Coût annuel de la déduction immédiate des équipements et ressources pour la financer}{fig:cout}

\begin{table}[!htbp]\centering\scriptsize
\caption{Les arbitrages de l'IS (Md€ de 2027 par an)}\label{tab:is_arb}
\begin{tabularx}{\linewidth}{@{}>{\raggedright\arraybackslash}p{3.4cm}>{\raggedright\arraybackslash}X>{\raggedright\arraybackslash}p{4.2cm}@{}}
\toprule
Choix & Ce qu'il change & Qui paie \\
\midrule
Déduction immédiate : 50~\% ou 100~\% & 50~\% : 7 à 7,8 la 1\up{re} année, gain de PIB divisé par deux ; 100~\% : 14 à 15,5, puis 1,4 à 2 & --- \\
Calendrier : 100~\% dès 2027, ou 50~\% jusqu'en 2032 puis 100~\% & dette maximale (intérêts à 3,25~\%) : 17 à 45 ou 2 à 32 ; remboursée en 10 à 28 ou 4 à 25~ans & --- \\
Financement par l'assiette & Fin du taux réduit (2,2 à 3,1), CIR apprécié au niveau du groupe (0,6), fin des exonérations zonées hors outre-mer (0,55), boîte à brevets à 17,5~\% (0,2 à 0,6) : 3,5 à 4,9 & petites sociétés ; grands groupes de recherche ; 1~065 utilisateurs de la boîte à brevets \\
Financement par le taux (28~\%) & 5,1 à 8,7 ; perte sèche +0,04 sur le capital et 0 à +3,8 selon le coût réel du transfert de bénéfices & toutes les sociétés, surtout les profits mobiles \\
Taux réduit maintenu & il manque 2,2 à 3,1 ; la dette de transition n'est jamais remboursée & --- \\
Taux réduit réservé aux sociétés opérationnelles & rendement non chiffrable (aucune ventilation publiée) & holdings patrimoniales, sociétés immobilières à l'IS \\
Excédent une fois la dette remboursée & rendu par une baisse du taux normal de 0,6 à 1,2 point & --- \\
Contribution exceptionnelle & s'éteint (référence) : $-$7 par rapport à 2026, $-$5 par rapport au PLF 2027 & environ 450 grands groupes \\
\bottomrule
\end{tabularx}
\end{table}

\textbf{Le taux ou l'assiette.} Une hausse du taux normal de 25 à 28~\% rapporterait 8,7~Md€ à bénéfices
inchangés et 5,1~Md€ si les bénéfices déclarés baissent de 1,5~\% par point, l'ordre de grandeur trouvé sur les
filiales de multinationales (Beer, de Mooij et Liu, 2020). L'objection n'est pas son montant, mais ce qu'elle
frappe : la déduction immédiate cesse d'imposer le rendement normal des investissements nouveaux, alors que le
taux porte d'abord sur les profits exceptionnels, qui sont aussi les plus mobiles. Un groupe choisit le pays
où il installe un projet très rentable selon le taux moyen effectif, et celui où il déclare ses bénéfices
selon le taux affiché. Si les bénéfices déplacés sont une perte réelle pour la France, la perte sèche de
28~\% atteint plusieurs milliards ; si ce ne sont que des écritures, elle est presque nulle (tableau~\ref{tab:is_arb}).
Le financement par l'assiette (3,5 à 4,9~Md€) dépasse le coût permanent de la déduction ; l'excédent rembourse
la transition, puis peut être rendu par une baisse du taux normal de 0,6 à 1,2 point.

\textbf{La transition} est la vraie difficulté. Une déduction totale dès 2027, financée par l'assiette,
creuserait une dette de 17 à 45~Md€, intérêts compris, remboursée en 10 à 28~ans. Un calendrier voté dès 2027,
50~\% jusqu'en 2032 puis 100~\%, limite la dette à 2 à 32~Md€, remboursée en 4 à 25~ans ; il évite de laisser le
second palier à l'appréciation des budgets futurs, ce qui reproduirait l'instabilité reprochée à l'IS.

\textbf{La fin du taux réduit : qui y perd.} Les sociétés dont le chiffre d'affaires est inférieur à 10~M€ et
le capital détenu à 75~\% par des personnes physiques paient 15~\% au lieu de 25~\% sur les 42\,500 premiers
euros de bénéfice ; depuis une décision du Conseil d'État de 2025, le chiffre d'affaires d'une société
membre d'un groupe s'apprécie au niveau du groupe. Le dispositif coûte environ 2,7~Md€ (bénéfices de 2022).
Pour une société au-dessus du plafond, il revient à verser chaque année 4\,250~€, quoi qu'elle fasse : la
somme ne dépend d'aucune décision, sauf celle de rester petite ou de se découper (on observe un pic de
sociétés juste sous le plafond). Mais beaucoup de bénéficiaires sont sous le plafond (environ 24\,000~€ de
bénéfice imposé à 15~\% en moyenne en 2018) : pour eux, la suppression relève l'impôt de chaque euro
supplémentaire, de deux tiers en tout (3\,000~€ pour 30\,000~€ de bénéfice), et renchérit tout leur capital.
Bernini et Treibich (2016) trouvent que l'introduction du taux réduit en 2001-2003 a accru de 29~\% le
capital des petites sociétés et leurs exportations, mais ralenti leur emploi (résumé seulement lu). La
déduction immédiate ne compense qu'une partie : pendant la phase à 50~\%, elle avance 2\,300~€ d'impôt la
première année pour 20\,000~€ d'équipements, et ne compense durablement la perte qu'au-delà de 220\,000~€
d'investissement annuel (110\,000~€ à 100~\%). Perdent le plus les sociétés qui investissent peu en
équipements : professions libérales, conseil, petites holdings et sociétés immobilières à l'IS ; aucune
statistique ne ventile la base imposée à 15~\% par type de société. Pour l'associé qui se fait distribuer le
bénéfice, le taux cumulé de l'IS et du PFU sur les 42\,500 premiers euros passe de 41,7 à 48,6~\%, celui de
toutes les autres sociétés.

\textbf{Ce que la référence écarte.} Un barème progressif de l'IS : la taille d'une entreprise ne dit rien de
la richesse de ses propriétaires, et les seuils multiplieraient les découpages. Un taux minimal national :
le minimum mondial de 15~\% s'applique déjà aux groupes de plus de 750~M€ de chiffre d'affaires (loi de
finances pour 2024, art.~33) ; pour les autres, un minimum assis sur l'impôt de l'année frapperait les
entreprises qui investissent beaucoup cette année-là. La contribution exceptionnelle s'éteint : surtaxe en
partie rétroactive et renouvelée d'année en année, elle fait payer les investissements passés et rend
incertain le rendement des investissements futurs.
